% Gedeeld door de syllabi van labo 3, 4 en 5. Pas dit hier aan, niet in de labobestanden.
\section{Software en lesmateriaal}

Elk labo is een \textbf{marimo-notebook} in de map \code{labs/} van de
cursusrepository. Je kunt het op twee manieren openen:

\begin{itemize}
    \item \textbf{Online:} surf naar \url{https://marimo.app}, kies
        \emph{New $\rightarrow$ Open from URL\dots} en plak de raw-URL van
        het notebook op GitHub. De data wordt dan van GitHub gelezen, dus
        je hoeft niets te installeren.
    \item \textbf{Lokaal:} installeer de pakketten uit
        \code{requirements.txt} en open het notebook met marimo. De data
        wordt dan uit de map \code{data/} gelezen.

\end{itemize}

Drie dingen die je moet weten over marimo:

\begin{itemize}
  \item Notebooks zijn \textbf{reactief}: als je een cel aanpast, worden
        alle cellen die ervan afhangen opnieuw uitgevoerd, in de volgorde
        van hun afhankelijkheden.
  \item \textbf{Twee cellen kunnen niet dezelfde variabele definiëren.}
        Geef elke stap een eigen, duidelijke naam (\code{tests\_raw},
        \code{tests\_clean}, \dots) of zet een reeks stappen samen in één
        cel.
  \item Elke oefening bevat een placeholder \code{...}. Een verborgen cel
        eronder controleert je antwoord: die kleurt \textbf{groen} als het
        juist is, en geeft anders een algemene hint over wat nog niet
        klopt, zonder het antwoord te verklappen.
\end{itemize}
